\documentclass{article}
\usepackage[letterpaper, margin=1in]{geometry}
\usepackage{fancyhdr}
\usepackage[utf8]{inputenc}
\usepackage{amsmath}
\usepackage{circuitikz}
\usepackage{pgfplots}
\usepackage{bm}
\usepackage{amsfonts}
\usepackage{graphicx}
\usepackage{float}
\usepackage{subcaption}
\usepackage[dvipsnames]{xcolor}
\usepackage{array}
\usepackage[colorlinks=true, allcolors=blue]{hyperref}
\title{6. Complex Power in AC Circuits \\  \large Introduction to Electric Circuits \\ Supplemental Instruction}
\author{Cooper McAllister}
\date{April 23 2026}
\begin{document}
\fancypagestyle{footerstyle} {
	\fancyhf{}
	\renewcommand{\headrulewidth}{0pt}
	\renewcommand{\footrulewidth}{0.4pt}
	\fancyfoot[L]{\textit{For personal study and students leading supplemental instruction, \\ with attribution given to the author. \\ Find more at coopermcallister.com}}
	\fancyfoot[R]{\thepage}
}
\pagestyle{footerstyle}
\maketitle
\thispagestyle{footerstyle}

\begin{center}
\textit{This document is not a substitute for a textbook or classroom instruction. It may contain errors and is intended to be used only to the extent that it is helpful to supplement students' learning experience.}
\end{center}

\section{Three Phase Line and Phase Voltages and Currents}
Three Phase circuits can be connected either in either wye (Y) or delta ($\Delta$) configuration. The relationship between line voltages and phase voltages, as well as between line currents and phase currents, can be shown to be as follows. These equations apply for magnitudes or RMS values.
\\In Y configuration:
$$V_l = \sqrt{3}V_p$$
$$I_l = I_p$$
In $\Delta$ configuration:
$$V_l = V_p$$
$$I_l = \sqrt{3}I_p$$
When looking at a circuit diagram, it is easy to see for each configuration which quantities are the same and which are different.


\newpage
\section*{References}
\begin{itemize}
\item M. Iordache. Class Lectures, Topic: ``Electric circuits.'' EEGR 2053, School of Engineering and Engineering Technology, LeTourneau University, 2025.
\item M. Iordache. 2024. EEGR 2053 Electric circuits [PowerPoint slides]. Available: \\ https://mviordache.name/EEGR2053/index.html
\item T. L. Floyd and D. M. Buchla, Principles of Electric Circuits: Conventional Current, 10th ed. Hoboken, NJ: Pearson Education, 2020.
\item O. Ortiz. 2025. ENGR 2053 Intro to electric circuits [PowerPoint slides].
\end{itemize}

	
\end{document}